\subsection{Query Model}
A query selector identifies a discovery class, leaf, and optional result index. The selector is a Q-sized value held
in the selected (:term:base.terms.rn_register:).

CPUID is unprivileged. For a (:term:base.terms.logical_processor:), CPUID results remain stable until
that (:term:base.terms.logical_processor:) is reset. A discovery result applies to execution on the (:term:base.terms.logical_processor:) that returned it.
Software associates cached discovery state with that (:term:base.terms.logical_processor:) and uses results obtained on each logical
processor before selecting its optional architectural behavior.

\begin{BedrockListedFormatDiagram}{CPUID Query Selector Format}
\BedrockFormatRowRange{selector[63:0]}{63}{0}{%
\BedrockFormatField{class}{32}
\BedrockFormatField{leaf}{16}
\BedrockFormatField{index}{16}
}
\end{BedrockListedFormatDiagram}

\BedrockTableCaption{CPUID Query Selector}
\begin{BedrockTabular}{@{}>{\raggedright\arraybackslash}p{0.65in}>{\raggedright\arraybackslash}p{0.85in}>{\raggedright\arraybackslash}p{3.90in}@{}}
\toprule
\rowcolor{BedrockHeaderFill}
\textbf{Bits} & \textbf{Field} & \textbf{Meaning}\\
\midrule
\texttt{0..15} & \texttt{index} & result index within the selected leaf\\
\texttt{16..31} & \texttt{leaf} & information leaf within the selected class\\
\texttt{32..63} & \texttt{class} & CPUID information class\\
\bottomrule
\end{BedrockTabular}

Index zero of every defined leaf is a common discovery (:term:base.terms.header:). Leaf-specific (:term:base.terms.payload:) begins at index one. An unknown
class, leaf, or index returns zero.

\begin{BedrockListedFormatDiagram}{CPUID Common Leaf Header Formats}
\BedrockFormatRowRange{class 0, leaf 0}{63}{0}{%
\BedrockFormatField{MAX\_CLASS}{32}
\BedrockFormatField{MAX\_LEAF}{16}
\BedrockFormatField{MAX\_INDEX}{16}
}
\BedrockFormatRowRange{other class, leaf 0}{63}{0}{%
\BedrockFormatReserved{(:term:base.terms.reserved:)}{32}
\BedrockFormatField{MAX\_LEAF}{16}
\BedrockFormatField{MAX\_INDEX}{16}
}
\BedrockFormatRowRange{leaf 1 or greater}{63}{0}{%
\BedrockFormatReserved{(:term:base.terms.reserved:)}{48}
\BedrockFormatField{MAX\_INDEX}{16}
}
\end{BedrockListedFormatDiagram}

\BedrockTableCaption{CPUID Common Leaf Header}
\begin{BedrockTabular}{@{}>{\raggedright\arraybackslash}p{1.85in}>{\raggedright\arraybackslash}p{1.15in}>{\raggedright\arraybackslash}p{1.05in}>{\raggedright\arraybackslash}p{1.05in}@{}}
\toprule
\rowcolor{BedrockHeaderFill}
\textbf{Query} & \textbf{Bits 63..32} & \textbf{Bits 31..16} & \textbf{Bits 15..0}\\
\midrule
class 0, leaf 0, index 0 & \texttt{MAX\_CLASS} & \texttt{MAX\_LEAF} & \texttt{MAX\_INDEX}\\
other class, leaf 0, index 0 & (:term:base.terms.reserved:), zero & \texttt{MAX\_LEAF} & \texttt{MAX\_INDEX}\\
any class, leaf 1 or greater, index 0 & (:term:base.terms.reserved:), zero & (:term:base.terms.reserved:), zero & \texttt{MAX\_INDEX}\\
\bottomrule
\end{BedrockTabular}

\texttt{MAX\_CLASS} is the greatest class recognized by the implementation, \texttt{MAX\_LEAF} is the greatest leaf
recognized within the selected class, and \texttt{MAX\_INDEX} is the greatest index recognized within the selected
leaf. Each (:term:base.terms.field:) gives the greatest recognized selector value. A leaf may be sparse below \texttt{MAX\_INDEX}. Software
ignores (:term:base.terms.reserved:) bits in every returned value.

CPUID bit (:term:base.terms.field|plural:) are feature predicates unless the leaf description says they encode a numeric value. Reserved
result bits are (:term:base.terms.ignored:) by software.

A set feature bit only permits software to use the corresponding architectural behavior; the instruction\textquotesingle{}s normal
privilege, operand, and availability checks still apply.

Numeric (:term:base.terms.field|plural:) use names that describe their unit. \texttt{SAVE\_AREA\_SIZE\_BYTES},
\texttt{OFFSET\_BYTES}, \texttt{MAX\_SIZE\_BYTES}, and \texttt{ALIGNMENT\_BYTES} are byte counts.
